\documentclass[11pt,a4paper]{article}
\usepackage[utf8]{inputenc}
\usepackage[T1]{fontenc}
\usepackage{geometry}
\usepackage{xcolor}
\usepackage{hyperref}
\usepackage{titlesec}

\geometry{margin=1in}
\definecolor{tumblue}{RGB}{0,101,189}

\hypersetup{
    colorlinks=true,
    linkcolor=tumblue,
    urlcolor=tumblue,
    pdftitle={Cover Letter - Ismail AISSAOUI},
}

\begin{document}

\noindent \textbf{Ismail AISSAOUI} \\
State Engineer in Artificial Intelligence (ESI-SBA) \\
Email: ismail.aissaoui@example.com \\
Phone: [Your Phone Number] \\
LinkedIn: [Your Link]

\bigskip

\begin{flushright}
    Professor Nils Thürey \\
    \textbf{Technical University of Munich (TUM)} \\
    Focus Group for Simulation and Digital Twin \\
    Munich, Germany \\
    \medskip
    May 5, 2026
\end{flushright}

\bigskip

\noindent \textbf{Subject: Application for Doctoral Research - Differentiable Physics \& Digital Twins}

\bigskip

Dear Professor Thürey,

I am writing to formally submit my application for a doctoral position within your Focus Group for Simulation and Digital Twin at the Technical University of Munich. As an AI Engineer with a deep interest in the synergy between deep learning and physical simulations, I have closely followed your work on differentiable physics and generative modeling for fluid dynamics.

Currently, I am finalizing my state engineering thesis at **Sonatrach**, where I have designed and implemented a **Generative Digital Twin** for the monitoring of gas turbine systems. My research addresses the challenge of making neural surrogates "physics-aware" by embedding conservation laws into the training process using **PINN** frameworks. I have found that this approach is critical for the reliability of digital twins in safety-critical industrial environments.

I am particularly inspired by your efforts to combine classical solvers with neural networks to create faster and more accurate simulation tools. My technical experience with **Graph Neural Networks (GNNs)** for modeling system topology and **Mamba-SSM** for high-performance sequential modeling allows me to approach physical systems from a multiscale perspective. I am eager to explore how these architectures can be used to develop **Differentiable Industrial Surrogates** that can be trained on a mix of simulation data and real-world sensor streams.

TUM is my first choice for a PhD because it is at the forefront of the Digital Twin revolution. I am a highly motivated researcher, proficient in Python and PyTorch, and I am eager to bring my industrial experience at Sonatrach to your research group.

I have attached my academic transcripts, my CV (noting ESI-SBA's H+ recognition), and a research proposal for your review. I would be honored to discuss how I can contribute to the state-of-the-art research in your laboratory.

Thank you for your time and consideration.

Sincerely,

\bigskip

\noindent Ismail AISSAOUI

\end{document}
